\section{Experiments}
\label{sec:experiments}
实验围绕三个问题组织：少量目标交互是否改善受约束任务表现；校准后的
分布迁移是否提供超出动作适应与模型控制的收益；这种收益能否推广到未见
机制和身体。本节给出拟采用的比较协议，正式结果表和结论留待独立评估。

\subsection{Tasks, Bodies, and Shared Control Interface}
源机器人为 Panda，固定臂目标为 xArm7，并以 UR5e 检验跨自由度迁移。
目标几何、关节边界、工具与控制能力对全部方法公开；未知控制响应和接触
参数通过声明的目标交互获得。身体变化与机制变化独立交叉，分别比较同身体
不同响应、不同身体相近响应，以及身体与响应同时变化。

四项任务使用任务空间技能 $U\in\R^{H\times13}$：
\begin{equation}
 U_j=\big[v_j^\top,\ \omega_j^\top,
 (\log k^{\rm trans}_j)^\top,
 (\log k^{\rm rot}_j)^\top,\ F^d_{n,j}\big]^\top.
 \label{eq:experimentalprimitive}
\end{equation}
其中 $v,\omega$ 为线速度与角速度，$k$ 为相对参考值的对角刚度，
$F_n^d$ 为法向力设定值。各身体通过自己的低层控制接口实现这些指令；
自由空间阶段屏蔽力指令。任务阈值、坐标系、控制周期、碰撞白名单、
终止规则和初态分布在开发阶段确定，并在正式评估前冻结。

\begin{enumerate}[label=\textbf{Task \arabic*.},leftmargin=2.1em]
\item \textbf{ObstacleTransfer.} 持带方向性的刚性工件绕过横梁与侧板，
穿过姿态受限的三维开口，到达指定 $SE(3)$ 位姿。源行为包含上绕与侧绕
路径；评估几何差异、执行延迟与负载变化对无碰撞完成、制动和跟踪的影响。
\item \textbf{GuidedContactAlignment.} 沿三维导面建立接触，通过滑移和旋转
到达定位区域。比较控制响应、摩擦与柔顺性变化，重点检验自由空间跟踪
相近时，是否仍需调整压靠、速度、刚度和接触序列。
\item \textbf{KeyedInsertionRecovery.} 从横向与角度偏移启动，将非轴对称
connector 插入匹配 socket；包含受限搜索、回撤、重对齐和再次推进。
重点评价整体完成、卡滞后的恢复及瞬时侧向力，而非仅评价平均接触力。
\item \textbf{MobileKeyedInsertionRecovery.} 保持 Task 3 的工件、夹具、
力限和恢复语义，改变身体支撑与接触耦合。检验固定臂技能在移动机械臂
和带臂腿式平台上的受约束实现。
\end{enumerate}

\begin{table}[ht]
\centering\small
\begin{tabularx}{\linewidth}{L{2.1cm}YY}
\toprule
任务 & 主要任务指标 & 机制与鲁棒性指标\\\midrule
Task 1 & 无碰撞完成、位姿误差、时间 & 全臂/持件间隙、跟踪误差、饱和\\
Task 2 & 限力定位、定位误差、时间 & 接触丢失、滑移/停滞、力峰值\\
Task 3 & 受约束插入、深度、时间 & 卡滞、条件恢复、重试、侧向力峰值\\
Task 4 & 插入且身体约束满足、总时间 & 支撑失败、滑移、身体偏移、WBC 饱和\\
\bottomrule
\end{tabularx}
\caption{四任务的比较重点。各任务均保留全部请求的失败与拒绝分母。}
\end{table}

\subsubsection{Task 4: Same Operation with Mobile Support}
\label{sec:task4}
Task 4 复用冻结的 Panda Task 3 prior。先比较固定与移动的同一机械臂，
再检验带臂腿式身体；真实轮地或足地动力学与平台约束分别定义。
平面运动代理只能用于接口诊断。跨移动支撑类型的结论需要多个有效支撑
类型的独立证据；人形平台取决于实际控制接口，不作为预设完成项。

同一身体上的全部高层方法共享 WBC、站位规则和可行性投影。
任务条件包含可用的基座姿态/速度、支撑接触、关节余量与控制器历史，
身体专属状态不直接塞入冻结的 Panda prior。末端响应对应允许为完成任务
所需的合法身体运动，身体稳定进入目标约束。
UMI-on-Legs 提供操作策略与全身控制分层的参照\cite{umilegs}；
改变 WBC 的系统比较与固定 WBC 的高层方法比较分别报告。

\subsection{Strong Baselines and Data Permissions}
\textbf{受控校准协议 A。} 全部方法获得相同源数据、公开任务目标、
已知身体模型、可用传感与低层控制接口。目标仅提供预算内短探针；不提供
完整任务成功演示或目标在线 RL。校准后的模拟训练只能使用由公开模型和
预算内记录构建的模型。隐藏目标环境仅提供计费探针及封存评估，
实际机制参数不向训练开放。

\textbf{冻结主表。} 主表固定为九行，并要求每行参加四个任务：本文方法、
matched task-only、FADA task-adapted、ACP task-adapted、
Human2Sim2Robot source-trajectory adaptation、DIAL-MPC task-adapted、
ordinary response MPC、public geometry/feedback 与 public source prior。
这一组合覆盖 2025/2026 年近期强方法、Karen Liu 与 Guanya Shi 参与的路线，
并保留三个不依赖论文适配细节的强公共控制\cite{human2sim,fada,acp,dialmpc}。

\begingroup\small
\begin{longtable}{L{3.2cm}L{2.2cm}L{2.7cm}L{4.5cm}}
\toprule
主表行 & 权限 & 身份 & 可解释的比较范围\\\midrule\endfirsthead
\toprule 主表行 & 权限 & 身份 & 可解释的比较范围\\\midrule\endhead
Ours gated transport & B-log & $C+H+GR$ & 与 task-only/FADA 做同权限任务比较；gate 机制另由六格消融识别。\\
Task-only & B-log & $C$ & 同容量、同数据、同训练流与同 checkpoint fallback 的首要对照。\\
FADA & B-log & Planner/IDM task adaptation & 保留公开 planner/IDM 核心；标为任务适配，不冒充原人形定量复现\cite{fada}。\\
ACP & C-supervised & virtual target/stiffness adaptation & 目标示范与力监督单列；不能与 B-log 作相同信息因果比较\cite{acp}。\\
Human2Sim2Robot & C-supervised-model & robot-source trajectory proxy & 原人类资产不可用时明确标为 source-trajectory 改编，不推广成原方法结论\cite{human2sim}。\\
DIAL-\newline MPC & B-log\newline model & annealed diffusion\newline online MPC & 保留在线 candidate rollout 与 receding horizon；模型调用与 deadline miss 进入资源列\cite{dialmpc}。\\
Ordinary response MPC & B-log\newline model & online Adam shooting & 与 DIAL-MPC 使用同一公共 response model 和 decoder/WBC。\\
Public geometry & no-target-learning & 固定几何/反馈 & 检验任务是否已可由强公共反馈解决。\\
Public source prior & no-target-learning & 冻结 source skill & 检验不做目标适配时的直接迁移与部署前拒绝。\\
\bottomrule
\end{longtable}
\endgroup

Play2Perfect、ASAP、SPI-Active 与 XMoP 保留为相关或后续扩展路线，
不在冻结九行中用不完整代理占据主表位置\cite{play2perfect,asap,spi,xmop}。
原生复现、完整系统比较与共同接口改编不能互相替代；每个 task-adapted 行
均报告 source、known-nominal 与目标三段漏斗。

\textbf{部署前拒绝。} 一个实现若耗尽预声明机会并在 source/known-nominal
资格门失败，仍覆盖该任务单元。其每个计划请求记
constrained success $=0$、constraint violation $=1$、capped normalized
time $=1$、pre-deployment rejection $=1$、target execution $=0$。
这里的 violation 是算法部署合同失败，不表示未执行目标轨迹发生了物理超限。
主表同时报告 rejection rate、target execution rate 与物理执行数。

DR/recurrent PPO、经典 Residual RL、通用校准 MPC、Mirage 和 ASID
保留为基础参照\cite{dr,ppo,rppo,residualrl,mjpc,mirage,asid}。
FORGE 是 RA-L 2025 接触参照，不列作已核验的近期顶会主会方法\cite{forge}。
同容量 task-only、deterministic transport、reweighting-only 和去响应/模式/
风险项属于内部消融。强对照可使用连续动作、技能 chunk、历史条件与合理
网络容量，不限制在有限映射候选库。

\textbf{共同权限的完整适应比较。} B-log 向本文方法和 FADA/ASAP 等对照
开放相同总时长、可测信息、reset 与执行约束下的完整自采 action/observation
轨迹，不含目标 reward 或专家示范，各方法自行选择采集策略。
另做共享同一记录缓冲的控制实验，区分采集与适应收益。
ASAP 所需状态对齐与额外测量权限单列；同组方法获得相同可用信息。
不使用 reward 不会使完整目标轨迹自动变成 A 短探针。
B-online 允许计费的目标任务交互与 reward；
C-supervised 比较向本文扩展版本提供相同可用示范并标明监督；信息不同的
原生系统另报性能--成本前沿，不作同预算排名。外部预训练单列。
D-oracle 仅作
具有目标真参数的封存诊断。需要大量随机目标数据的 latent alignment、
目标演示训练的 diffusion/latent-action policy，以及目标在线 residual RL
分别报告数据曲线\cite{lsa,lad}。充分目标训练的策略与真值规划器不并入
少量校准主表，也不被解释为已证明的性能上界。

\textbf{接口与训练公平性。} 受控比较共享传感、可用历史、公开任务 reward、
控制周期、命令边界和安全判据。需要特权 critic 的方法只能读取其训练模型的
状态，隐藏目标参数不得进入训练。原生实现的复现检查与统一任务接口的改编
结果分别记录，不能把本项目同名策略直接视为原论文复现。

\subsection{Budget, Generalization, and Statistical Evaluation}
目标访问按秒数、transitions、reset 和数据类型报告。
短探针子协议可采用 $B\in\{0,5,10,30\}$，每次时长与 fit/validation
分配在正式测试前冻结。需要约两分钟动作/观测记录的 FADA 等方法，
另包含原生时长的数据点\cite{fada}，不能用不合适的极短缓冲代表完整性能。
目标轨迹即使不含 reward，也须声明是否为完整任务执行；原生示范另列。
接口调整、验证与失败尝试全部计费，主动方法受相同总访问预算限制。

训练模拟 transitions、离线更新、调参次数、GPU 时间和部署在线模型调用
另列。比较同模拟量的学习曲线与固定时延下的表现；不同优化器的更新次数
不直接视为相同计算量。MPC 超时须按预先固定的执行规则记入结果。

测试分别留出初态、机制组合与身体。UR5e 的两项协议分开：
冻结 transport 仅拟合 context；以及允许在校准模型内继续更新策略的
few-shot adaptation。已知目标形态上的训练不等于未见身体的冻结泛化。
仿真任务结果与实机结果单列；平台与传感接口未确定时不预设实机结论。

正式评估冻结为八个独立训练/校准实例、两个目标机制组和配对请求。
Task 1--3 每实例为 $10$ 个不同物理初态 $\times2$ 个请求模式
$\times2$ 个机制组，共 $40$ 个请求；Task 4 再交叉两个 body/support cells，
共 $80$ 个请求。确定性公共控制只保留一个方法身份，在同一 request package
上执行，不伪造八个训练重复。所有方法共享相同 task/mechanism/body/initial/mode
键，失败不替换。完整 post-reset 物理与控制器状态 hash 才定义物理初态；
episode 标签或相同 deterministic reset 上的 policy seed 不构成新初态。

八实例由开发数据上的预声明配对 discordance power simulation 冻结。
中等 discordance 场景达到目标 power；高训练异质性仍可能不足，作为限制报告，
不能在查看最终结果后追加实例。推断以训练/校准实例为最高随机层，再对
mechanism/body/initial request block 做 $10{,}000$ 次层级配对 bootstrap。
每个研究问题预声明一个主对比，RQ family 内使用 Holm 修正；报告效应量、
区间和原始分子/分母，不给出无权跨任务总排名。全部支持失败、验证拒绝、
超力、超时和身体失败进入整体分母。

\subsection{Aligned Six-Cell Ablation}
六格消融与主方法使用相同 warm checkpoint、数据流、优化器、更新数、
held-out 模型代理和 task-only $+0.05$ fallback：
\begin{center}
\begin{tabular}{lll}
\toprule
单元 & 目标 & 识别作用\\\midrule
frozen-warm & 不继续训练 & continued training 控制\\
task-only & $C$ & task-first 基准\\
matching-no-band & $C+R$ & response matching 主效应\\
band-only & $C+H$ & additive task band 主效应\\
ungated-full & $C+H+R$ & 无 gate 的 $2\times2$ 完整角\\
gated-full & $C+H+GR$ & selective gate\\
\bottomrule
\end{tabular}
\end{center}
$C,C+R,C+H,C+H+R$ 是严格的无 gate $2\times2$；只用
gated-full$-$ungated-full 识别 gate。六格都参加四任务记账；若整个任务族在
冻结资格门失败，则六格均按部署前拒绝进入分母，不训练或查看目标结果。
历史上用 hinge $H$ 替换 $C$ 的五格工件与当前方法身份不一致，不进入推断。

\subsection{Performance Evidence and Contribution Tests}
主要评价无碰撞且满足力/身体约束的任务成功、失败风险、完成时间，以及
获得这些表现所需的目标数据与计算成本。dense reward 用于训练；统一的任务
成功判据用于评估，累计 reward 不能跨任务代替成功与安全表现。

\begin{enumerate}
\item \textbf{交互效率。} 在同等目标访问权限下，比较成功率--预算曲线，
以及达到预声明成功水平所需的目标交互量。
\item \textbf{接触与恢复。} 在相同约束下比较接触定位、卡滞后的整体恢复、
力峰值与失效概率，并检验未见机制组合。
\item \textbf{任务--行为权衡。} 首先控制任务成本，再比较有任务意义的模式
覆盖、响应和条件对应。无模式请求时允许 baseline 选择更可靠的模式；
有明确模式请求时，向全部方法提供相同条件，并分别评价按请求完成和自主完成。
\item \textbf{摊销执行。} 在与校准模型 MPC 相当的任务表现下，检验 transport
是否降低端到端在线时延；WBC、状态处理与模型计算均计入。
\end{enumerate}
源--目标响应距离下降、额外噪声或更大模型本身不构成性能贡献。
方法是否成立取决于上述比较，不能预设所有任务或指标均优于现有方法。
